\documentclass[11pt,a4paper]{article}
\usepackage[T1]{fontenc}
\usepackage[utf8]{inputenc}
\usepackage{lmodern}
\usepackage[margin=2.15cm]{geometry}
\usepackage{setspace}
\setstretch{1.08}
\setlength{\parskip}{0.25em}
\setlength{\parindent}{1em}
\usepackage{amsmath,amssymb}
\usepackage{booktabs}
\usepackage[hidelinks]{hyperref}

\title{Schedule-Aware Fan Control for\\Brewery Wastewater Cleaning}
\author{Short paper on innovation 1 of the thesis}
\date{}

\begin{document}
\maketitle
\thispagestyle{empty}

\begin{abstract}
\noindent
Aeration blowers are the main electrical load when a brewery cleans process water before discharge. The hard disturbances are not a daily sewage curve. They are clean-in-place recipes the production clock already names, including detergent pulses that cut oxygen transfer. This paper specifies a schedule-aware fan law: a one-minute inversion of the oxygen balance that previews the recipe, stages identical centrifugal blowers off their surge line, and falls back to a dissolved-oxygen loop if the optimizer stops. On a representative week the law uses 9656\,kWh of blower electricity, 11.6\,percent less than a proportional-integral loop parked at 2.5\,mg/L and 69\,percent less than constant speed sized for the detergent hour, with no hours above a study limit of 125\,mg/L chemical oxygen demand. A companion paper treats the live transfer factor. Here that factor is an input.
\end{abstract}

\section{Introduction}

Brewery wastewater is warm, strong, and scheduled. Equalization, anaerobic pretreatment, and aerobic polishing are the usual train when the water is cleaned and discharged rather than reused \cite{simate2011,fillaudeau2006,ashraf2021}. After anaerobic pretreatment the polishing tank is still an oxygen-transfer problem, and the fans that solve it are typically more than half of a biological plant's electricity \cite{longo2016,amand2013}. Published predictive aeration controllers act on municipal benchmarks whose input is a smooth airflow or a transfer coefficient, not a staged blower, and whose disturbance is a diurnal file \cite{holenda2008,shen2008,zeng2015,cristea2011}. They do not read a batch recipe.

The contribution is the disturbance and the actuator together. The brewhouse and clean-in-place calendar is encoded as a preview. The command is header airflow and the number of online machines, subject to a surge margin. Inner loops and interlocks stay in the programmable controller \cite{qin2003,obrien2011}. The numerical week below uses the deployed one-step law, not a mixed-integer solver. The horizon problem is the specification that law reduces to while the machines are identical.

\section{Plant and preview}

The aerobic volume is \(V=900\,\mathrm{m^3}\). Three identical variable-speed blowers share a header. Delivered airflow \(q\) is in normal cubic metres per hour, with a mixing floor of 900 and a header cap of 9000. Header gauge pressure is mostly static submergence,
\begin{equation}
\label{eq:dp}
\Delta p = 52 + 0.55\,(q/1000)^{2}
\quad[\mathrm{kPa}].
\end{equation}
Electrical power is polytropic compression of the delivered mass flow at inlet \(20^{\circ}\mathrm{C}\), with combined efficiency \(0.70\times 0.95\times 0.97\). Because \(\Delta p\) hardly grows with flow, power is much closer to linear in \(q\) than a fan cube law. The cube law remains appropriate only for the small exhaust fans on equalization and flotation, which are metered and are not decision variables here.

Oxygen transfer in the tank is
\begin{equation}
\label{eq:kla}
K_La = \alpha\,k_{\mathrm{ref}}\left(\frac{q}{q_{\mathrm{ref}}}\right)^{0.82}\theta^{\,T-20},
\end{equation}
with \(k_{\mathrm{ref}}=6.5\,\mathrm{h^{-1}}\) at \(q_{\mathrm{ref}}=2500\), \(\theta=1.024\), and \(\alpha\) the process-to-clean-water factor. Dissolved oxygen \(C\) obeys
\begin{equation}
\label{eq:do}
\dot C = K_La\big(\beta C^{*}(T)-C\big) - \mathrm{OUR},
\end{equation}
with \(\beta=0.95\). The calendar maps each recipe class to flow, chemical oxygen demand, temperature, a detergent flag, and a dissolved-oxygen set-point: 1.8\,mg/L in base production, 2.0 in the brewhouse peak and packaging clean-in-place, 2.2 in fermentation clean-in-place, and 1.3 on the weekend. The preview is applied 0.75\,h ahead of a boundary so the blower lag (0.10\,h) and the probe lag (0.08\,h) sit inside the lead.

\section{Deployed fan law}

Each minute the law reads the current regime \(\rho\) and the regime \(\rho^{+}\) at \(t+0.75\,\mathrm{h}\), and sets
\begin{equation}
C_{\mathrm{sp}}=\max\big\{C_{\mathrm{sp}}(\rho),\,C_{\mathrm{sp}}(\rho^{+})\big\}.
\end{equation}
Mass load \(L=Q\cdot\mathrm{COD}\) supplies a scale
\begin{equation}
\sigma=\mathrm{clip}(L^{+}/L,\,1,\,1.35).
\end{equation}
The scale is not allowed below one in the integrated week. That omission causes a short unload dip at the weekend and is a known defect, not a hidden one. Airflow is the inverse of \eqref{eq:kla} at uptake \(\sigma\,\mathrm{OUR}\) and driving force \(\beta C^{*}-C_{\mathrm{sp}}\):
\begin{equation}
\label{eq:inv}
q^{\star}
=
q_{\mathrm{ref}}
\left(
\frac{\sigma\,\mathrm{OUR}}
{\alpha\,k_{\mathrm{ref}}\,\theta^{\,T-20}\,(\beta C^{*}-C_{\mathrm{sp}})}
\right)^{1/0.82}.
\end{equation}
A clamped bias \(b\), with \(|b|\le 0.25\) and gain \(0.55\,\mathrm{h^{-1}}\) on \(C_{\mathrm{sp}}-C\), removes constant offset \cite{pannocchia2003}. The command is \(\mathrm{clip}((1+b)q^{\star},900,9000)\). Staging is one machine below \(2600\,\mathrm{Nm^{3}/h}\), two up to 5200, and three above. Surge flow at rated speed is \(0.22\times 2800=616\,\mathrm{Nm^{3}/h}\) per machine; the margin used is 1.15 times that, about \(708\,\mathrm{Nm^{3}/h}\). One machine at the mixing floor of 900 is legal. Two machines sharing 900 are not, so the optimizer must be allowed to turn machines off.

At a quiet base point (\(\mathrm{OUR}\approx 18\,\mathrm{mg\,L^{-1}\,h^{-1}}\), \(C_{\mathrm{sp}}=1.8\), \(T=28^{\circ}\mathrm{C}\), \(\alpha=0.60\)) the inversion returns about \(1550\,\mathrm{Nm^{3}/h}\) and about 33\,kW. One machine. At fermentation conditions (\(\mathrm{OUR}\approx 30\), \(C_{\mathrm{sp}}=2.2\), \(T=33^{\circ}\mathrm{C}\), \(\alpha=0.40\)) it returns about \(5250\,\mathrm{Nm^{3}/h}\) and about 133\,kW, on the two-to-three machine boundary. The same inversion with the quiet-week factor \(\alpha=0.55\) asks for only about \(3570\,\mathrm{Nm^{3}/h}\). The missing factor is \((0.55/0.40)^{1/0.82}\approx 1.47\). Innovation 1 can apply that factor only if innovation 2, or a recipe table, supplies \(\alpha\). In this paper \(\alpha\) is assumed known through the 0.30\,h filter used in the thesis estimate.

If the edge computer does not publish, the programmable controller holds a proportional-integral loop at 2.5\,mg/L. Low dissolved oxygen, a drive fault, a shut outlet valve, or header over-pressure override the set-point. The optimizer proposes. Level 1 disposes.

The economic problem this law approximates minimizes, over a horizon, blower energy plus hinge penalties on \(C\) below a floor and on chemical oxygen demand above 125\,mg/L, plus a penalty on each machine start. With identical machines the integer part is the rule above. With unequal machines the same slacks become three convex programs, one per machine count \cite{bemporad1999,rawlings2017}.

\section{Estimate}

The week is 96\,h base, 30\,h brewhouse, 18\,h fermentation clean-in-place, 12\,h packaging clean-in-place, and 12\,h weekend. Aerobic influent chemical oxygen demand runs from 160 to 460\,mg/L depending on regime. Four loops share the plant: the schedule-aware law; a predictive inversion with \(\alpha\) locked at 0.55 and no load preview; the 2.5\,mg/L proportional-integral fallback; and constant speed at 6552\,Nm\(^{3}\)/h, 1.15 times a detergent design point. Constant speed draws 185\,kW for the whole week because \eqref{eq:dp} at that flow is the design point of the power map.

\begin{table}[ht]
\centering
\caption{Blower electricity, representative week. COD violation hours are time above 125\,mg/L.}
\label{tab:week}
\small
\begin{tabular}{@{}lrrrr@{}}
\toprule
Loop & kWh & Mean kW & COD viol.\ (h) & Min DO in ferm.\ CIP \\
\midrule
Schedule-aware & 9656 & 57.5 & 0 & 1.63 \\
Fixed \(\alpha\), no preview & 9620 & 57.3 & 0 & 0.90 \\
PI at 2.5\,mg/L & 10920 & 65.0 & 0 & 1.87 \\
Constant speed & 31151 & 185.4 & 0 & 2.90 \\
\bottomrule
\end{tabular}
\end{table}

Table~\ref{tab:week} is the result. Every loop meets the study chemical-oxygen-demand limit, so the comparison is not a compliance contest. Against constant speed the schedule-aware law saves \(31151-9656=21495\,\mathrm{kWh}\) (69\,percent). Against the proportional-integral habit it saves \(1264\,\mathrm{kWh}\) (11.6\,percent). About half of that 11.6\,percent is the 96\,h base regime, where the set-point is 1.8 rather than 2.5 (3469\,kWh against 4102\,kWh). During fermentation clean-in-place the two use almost the same energy (2441 against 2460\,kWh). The fixed-\(\alpha\) inversion matches the schedule-aware energy and lets dissolved oxygen fall to 0.90\,mg/L in that window, under the 1.5\,mg/L operational floor. Preview and a live \(\alpha\) are what hold the floor at 1.63\,mg/L. They are not a second large energy cut once the set-point table exists.

The schedule-aware minimum elsewhere is not perfect. Entering the brewhouse peak, dissolved oxygen touches 1.39\,mg/L. Entering the weekend it touches 0.86\,mg/L while the weekend mean is 4.46\,mg/L, because the mixing floor over-aerates after the step. Allowing \(\sigma<1\) when the previewed load falls removes the second dip. It was not re-run. Exhaust fans on the same calendar add about 1300\,kWh and do not change the ranking.

\section{What to install first}

The percentages are not an instruction to deploy the edge computer on day one. Constant speed at the detergent-hour valve is the policy that wastes the 21495\,kWh. Moving that header onto the proportional-integral fallback already in the programmable controller takes most of the saving and does not require a calendar feed. The set-point table in this paper is the next increment, and it can run with \(\alpha\) taken from the recipe table rather than from a sensor. Fermentation clean-in-place then uses 0.40 instead of a quiet-week 0.55 even before any network is trained. Only after that week is logged against the fallback should a live \(\alpha\) close the inversion. The companion paper is that third step. Reversing the order puts a model in the loop while the plant still has no tested fallback \cite{obrien2011}.

A header that does not match \eqref{eq:dp} falsifies the kilowatt-hours and usually not the ranking against constant speed, because constant speed is the same wrong map held at one flow. The ranking against the proportional-integral loop is only twelve percent and is more fragile. The acceptance test is therefore three statements on a production week, not a copied percentage: blower energy below the fallback on the same calendar, hours above the site permit not higher than the fallback, and the dissolved-oxygen minimum during each detergent recipe above the floor the site has written down. In this study that floor is 1.5\,mg/L. The 11.6 and 69\,percent figures are the output of one integrator. The three statements are the result that travels.

\section{Conclusion}

A brewery fan controller should read the production clock and the blower map, not only a dissolved-oxygen error. On the representative polishing week, schedule-aware set-points and a surge-safe staging rule cut about twelve percent of blower electricity relative to a competent loop at 2.5\,mg/L, and about two thirds relative to constant speed sized for the detergent hour, without adding hours above 125\,mg/L chemical oxygen demand. The same week shows the limit of the claim: energy does not fall much further once a predictive law already uses a sensible set-point, unless the transfer factor is wrong, in which case the loss is dissolved-oxygen margin rather than kilowatt-hours. That margin is the subject of the companion paper. The figures are an integration of a stated model, not a trial on a named factory.

\begin{thebibliography}{99}
\bibitem{simate2011} G.~S. Simate et al. The treatment of brewery wastewater for reuse: State of the art. \textit{Desalination}, 273:235--247, 2011.
\bibitem{fillaudeau2006} L.~Fillaudeau, P.~Blanpain-Avet, and G.~Daufin. Water, wastewater and waste management in brewing industries. \textit{J. Cleaner Production}, 14:463--471, 2006.
\bibitem{ashraf2021} A.~Ashraf, R.~Ramamurthy, and E.~R. Rene. Wastewater treatment and resource recovery technologies in the brewery industry. \textit{Sustainable Energy Technologies and Assessments}, 47:101432, 2021.
\bibitem{longo2016} S.~Longo et al. Monitoring and diagnosis of energy consumption in wastewater treatment plants. \textit{Applied Energy}, 179:1251--1268, 2016.
\bibitem{amand2013} L.~{\AA}mand, G.~Olsson, and B.~Carlsson. Aeration control --- a review. \textit{Water Science and Technology}, 67:2374--2397, 2013.
\bibitem{holenda2008} B.~Holenda, E.~Domokos, {\'A}.~R{\'e}dey, and J.~Fazakas. Dissolved oxygen control of the activated sludge wastewater treatment process using model predictive control. \textit{Computers and Chemical Engineering}, 32:1270--1278, 2008.
\bibitem{shen2008} W.~Shen, X.~Chen, and J.-P. Corriou. Application of model predictive control to the BSM1 benchmark. \textit{Computers and Chemical Engineering}, 32:2849--2856, 2008.
\bibitem{zeng2015} J.~Zeng and J.~Liu. Economic model predictive control of wastewater treatment processes. \textit{Ind. Eng. Chem. Res.}, 54:5710--5721, 2015.
\bibitem{cristea2011} S.~Cristea, C.~de~Prada, D.~Sarabia, and G.~Guti{\'e}rrez. Aeration control of a wastewater treatment plant using hybrid NMPC. \textit{Computers and Chemical Engineering}, 35:638--650, 2011.
\bibitem{qin2003} S.~J. Qin and T.~A. Badgwell. A survey of industrial model predictive control technology. \textit{Control Engineering Practice}, 11:733--764, 2003.
\bibitem{obrien2011} M.~O'Brien, J.~Mack, B.~Lennox, D.~Lovett, and A.~Wall. Model predictive control of an activated sludge process: A case study. \textit{Control Engineering Practice}, 19:54--61, 2011.
\bibitem{pannocchia2003} G.~Pannocchia and J.~B. Rawlings. Disturbance models for offset-free model-predictive control. \textit{AIChE Journal}, 49:426--437, 2003.
\bibitem{bemporad1999} A.~Bemporad and M.~Morari. Control of systems integrating logic, dynamics, and constraints. \textit{Automatica}, 35:407--427, 1999.
\bibitem{rawlings2017} J.~B. Rawlings, D.~Q. Mayne, and M.~M. Diehl. \textit{Model Predictive Control: Theory, Computation, and Design}. Nob Hill, 2nd ed., 2017.
\bibitem{rosso2006} D.~Rosso and M.~K. Stenstrom. Surfactant effects on \(\alpha\)-factors in aeration systems. \textit{Water Research}, 40:1397--1404, 2006.
\end{thebibliography}

\end{document}
